\subsection{Các bảng trọng tâm}
\label{subsec:db_bang_trong_tam}

Bảng~\ref{tab:db_core_tables} chỉ nêu các nhóm bảng cần thiết để hiểu cách hệ thống lưu dữ liệu; chi tiết cột đã được lược bớt trong ERD.

\begin{table}[H]
\centering
\small
\begin{tabularx}{\textwidth}{|L{3.2cm}|L{4.1cm}|X|}
\hline
\textbf{Miền} & \textbf{Bảng đại diện} & \textbf{Vai trò} \\
\hline
Tài khoản và cảnh báo & \code{users}, \code{alerts}, \code{alert\_history} & Quản lý quyền sở hữu, cấu hình theo dõi và lịch sử thông báo. \\
\hline
Token và thị trường & \code{token\_meta}, \code{token\_market\_data}, các bảng chart/pool/holder & Lưu metadata, snapshot và chuỗi thời gian có nhịp cập nhật riêng. \\
\hline
Phân tích ví & \code{wallet\_overview\_cache}, \code{wallet\_analyses}, \code{wallet\_balance\_history} & Cung cấp tổng quan, PnL, win rate và biến động tài sản theo kỳ. \\
\hline
Hoạt động ví & Các bảng transfer/swap và \code{wallet\_enhanced\_*} & Phục vụ danh sách hoạt động, phân trang và chi tiết giao dịch on-chain. \\
\hline
\end{tabularx}
\caption{Các nhóm bảng trọng tâm của Yoca}
\label{tab:db_core_tables}
\end{table}

\code{wallet\_analyses} là read model chính của phần phân tích ví. Khóa \code{wallet\_address, period} cho phép lưu riêng kết quả theo từng kỳ; các chỉ số tổng hợp và phân rã theo ngày được giữ cùng một bản phân tích để giao diện nhận được dữ liệu nhất quán trong một lần đọc. Chi tiết ở cấp token và chuỗi số dư được lưu riêng vì có nguồn dữ liệu và nhịp cập nhật khác.

Index được đặt theo mẫu truy vấn. Lịch sử swap và transfer dùng address cùng timestamp để lấy hoạt động của một ví theo thứ tự thời gian. Các bảng enhanced transaction dùng address, signature và vị trí phần tử để chống ghi trùng khi đồng bộ lại. Dữ liệu thuộc người dùng được index theo user ID để mọi truy vấn profile, cảnh báo, chat và thanh toán luôn bắt đầu từ chủ sở hữu đã xác thực.
